\documentclass[10pt]{article}

\usepackage{amsmath}
\usepackage{amssymb}
\usepackage{amsthm}
\usepackage{ mathrsfs }
\usepackage{ mathtools}
\usepackage{enumerate}
\usepackage{bbm}
\usepackage{lipsum}
\usepackage{fancyhdr}
\usepackage{tikz-cd} 
\usepackage{adjustbox} 
\usetikzlibrary{arrows}
\newcommand{\midarrow}{\tikz \draw[-triangle 90] (0,0) -- +(.1,0);}
\tikzset{commutative diagrams/.cd,
mysymbol/.style={start anchor=center,end anchor=center,draw=none}
}
\newcommand\comsymb[2][\circlearrowleft]{%
  \arrow[mysymbol]{#2}[description]{#1}}

\newtheorem{theorem}{Theorem} [section] 
\newtheorem{proposition}{Proposition}[section] 
\newtheorem{definition}{Definition}[section] 
\newtheorem{lemma}{Lemma}[section] 
\newtheorem{notation}{Notation}[section] 
\newtheorem{remark}{Remark}[section] 
\newtheorem{corollary}{Corollary} [section] 
\newtheorem{terminology}{Terminology}[section] 
\newtheorem{fact}{Fact}[section] 
\newtheorem{example}{Example}[section] 
\newtheorem{claim}{Claim}
\newenvironment{claimproof}[1]{\par\noindent\underline{Proof:}\space#1}{\hfill $\blacksquare$}

\DeclareMathOperator{\tmod}{mod}
\DeclareMathOperator{\triv}{triv}
\DeclareMathOperator{\ind}{ind}
\DeclareMathOperator{\straw}{s.t.}
\DeclareMathOperator{\ev}{ev}
\DeclareMathOperator{\pr}{pr}
\DeclareMathOperator{\Aut}{Aut}
\DeclareMathOperator{\Map}{Map}
\DeclareMathOperator{\Stab}{Stab}
\DeclareMathOperator{\Ind}{Ind}
\DeclareMathOperator{\GL}{GL}
\DeclareMathOperator{\SL}{SL}
\DeclareMathOperator{\SO}{SO}
\DeclareMathOperator{\Sp}{Sp}
\DeclareMathOperator{\esssup}{esssup}
\DeclareMathOperator{\abs}{abs}
\DeclareMathOperator{\rel}{rel}
\DeclareMathOperator{\diam}{diam}
\DeclareMathOperator{\rank}{rank}
\DeclareMathOperator{\Hom}{Hom}
\DeclareMathOperator{\Homk}{Hom_{k-alg}}
\DeclareMathOperator{\Ker}{Ker}
\DeclareMathOperator{\Image}{Im}
\DeclareMathOperator{\Dom}{Dom}
\DeclareMathOperator{\grad}{grad}
\DeclareMathOperator{\divergence}{div}
\DeclareMathOperator{\rk}{rk}
\DeclareMathOperator{\Span}{Span}
\DeclareMathOperator{\mSpec}{mSpec}
\DeclareMathOperator{\MaxSpec}{MaxSpec}
\DeclareMathOperator{\interior}{int}
\DeclareMathOperator{\supp}{supp}
\DeclareMathOperator{\id}{id}
\DeclareMathOperator{\sgn}{sgn}
\DeclareMathOperator{\Mat}{Mat}
\DeclareMathOperator{\PD}{PD}
\DeclareMathOperator{\ext}{ext}
\DeclareMathOperator{\vol}{vol}
\DeclareMathOperator{\inte}{int}
\DeclareMathOperator{\diverg}{div}
\DeclareMathOperator{\curl}{curl}
\DeclareMathOperator{\image}{im}
\DeclareMathOperator{\dR}{dR}
\DeclareMathOperator{\Ob}{Ob}
\DeclareMathOperator{\Mnfd}{Mnfd}
\DeclareMathOperator{\OpEmb}{OpEmb}
\DeclareMathOperator{\Spec}{Spec}
\DeclareMathOperator{\Spa}{Spa}
\DeclareMathOperator{\Sper}{Sper}
\DeclareMathOperator{\op}{op}
\DeclareMathOperator{\con}{con}
\DeclareMathOperator{\Gal}{Gal}
%\newcommand*{\name}[\num_arguments][default values]{{\color{#1}\Large #2}}
\newcommand{\defeq}{\vcentcolon=}
\newcommand{\Ouv}{\mathcal{O}}
\newcommand{\norm}[1]{\Vert #1 \Vert}
\newcommand{\opNorm}[2]{\norm{#1}_{#2\to#2}}
\newcommand{\normL}[3]{\norm{#1}_{L^{#2}(#3)}}
\newcommand{\N}[0]{\mathbb{N}}
\newcommand{\m}[0]{\mathfrak{m}}
\newcommand{\R}[0]{\mathbb{R}}
\newcommand{\Z}[0]{\mathbb{Z}}
\newcommand{\C}[0]{\mathbb{C}}
\newcommand{\Q}[0]{\mathbb{Q}}
\newcommand{\Qc}[0]{\mathfrak{Q}_C}
\newcommand{\F}[0]{\mathbb{F}}
\newcommand{\G}[0]{\mathbb{G}}
\newcommand{\A}[0]{\mathbb{A}}
\newcommand{\T}[0]{\mathbb{T}}
\newcommand{\Para}[0]{\mathbb{P}}
\newcommand{\M}[0]{\mathcal{M}}
\newcommand{\maniN}[0]{\mathcal{N}}
\newcommand{\cinf}[0]{\mathcal{C}^\infty}
\newcommand{\torus}[0]{\mathbb{T}}
\newcommand{\sep}[0]{\:|\:}
\newcommand{\pd}[2]{{\frac {\partial #1} {\partial #2}}}
\newcommand{\compinc}[0]{\subset \subset}
\newcommand{\sH}[0]{\mathscr{H}}
\newcommand{\ab}[1]{\vert #1 \vert}
\newcommand{\st}[0]{\:\straw\:}
\newcommand{\g}[0]{\mathfrak{g}}
\newcommand{\p}[0]{\mathfrak{p}}
\newcommand{\U}[0]{\mathfrak{U}}
\newcommand{\fib}[1]{%
  \mathbin{\mathop{\times}\limits_{#1}}%
}
\newcommand{\tens}[1]{%
  \mathbin{\mathop{\otimes}\displaylimits_{#1}}%
}

\title{Rigid Analytic Geometry Week 5 Solution}
\author{So Murata}
\date{SoSe 26, University of Bonn}


\begin{document}
\maketitle

\par\textbf{Problem 1}\\
If $x=0$ or $y=0$, then trivially the statement holds, thus we assume $x,y\in K^\times$. Then $xy^{-1}\in R$ or $yx^{-1}\in R$. This proves the statement.
\par\textbf{Problem 2}\\
$\m_R=R\backslash R^\times$ is true by definition. Let $x\in R,y\in\m_R$, if $xy\in R^\times$ then $y^{-1} = y^{-1}x^{-1}x \in R$ thus $xy\in\m_R$. 
Note that $x^{-1}(x+y)\in\m_R$, then by the previous argument $x+y\in\m_R$. Thus without loss of generality, we assume $x|y$ in $R$. This is justified as $R$ is a valuation ring. Then $y=rx$ for some $r\in R$. If $x+y$ is invertible, then
\begin{equation*}
    x(r+1)
\end{equation*}
is invertible. This shows both $x,r+1$ are invertible as $x,r+1\in R$. We derived a contradiction and conclude $\m_R$ is closed under addition.
\par\textbf{Problem 3}\\
The direction $\Leftarrow$ is clear as any proper ideal is contained in $R\backslash R^\times$. The other direction follows from that any non-unital element is contained in some maximal ideal thus $R\backslash R^\times\subseteq\m$ where $\m$ is the unique maximal ideal. As $\m$ is proper, this is an equality.
\par\textbf{Problem 4}\\
Suppose for $x\in K$ be integral over $R$ and $x\not\in R$. Then $x^{-1}\in R$ and there are $a_1,\cdots,a_n\in R$ such that 
\begin{equation*}
    x^n+a_1x^{n-1}+\cdots+a_n = 0.
\end{equation*}
Then,
\begin{equation*}
    x = -(a_1+a_2x^{-1}+\cdots+a_nx^{n-1}).
\end{equation*}
All the elements in the right hand side are in $R$ thus $x\in R$. Thus $x\in R$.
\par\textbf{Problem 5}\\
Let $I,J$ be ideals of $R$. Suppose there are elements $a\in I,b\in J$ such that $a\not\in J,b\not\in I$. But either $a\in bR$ or $b\in aR$ thus this is not possible. $I,J$ are linearly ordered so is the set of whole ideals.
\par\textbf{Problem 6}\\
Let $\{a_1,\cdots,a_n\}\subseteq R$ then this is linearly ordered by inclusion $(a)\subseteq (b)$ for some $a_i,a_j=a,b$. Let $a$ be the maximal element, then,
\begin{equation*}
    (a_1,\cdots,a_n) = (a).
\end{equation*}
\par\textbf{Problem 7}\\
Prime ideals are equal to their own radicals and do not contain $1$ as they are proper. If $I$ is a proper ideal, then $1\not\in I$ thus $1\not\in\sqrt{I}$ otherwise it contains an unit. Let $ab\in\sqrt{I}$. Without loss of generality, we assume $b=ra$ for some $r\in R$. This is justified by the linearlity of ordering in ideals. Then,
\begin{equation*}
    b^2 = b\cdot ra\in \sqrt{I}.
\end{equation*}
Thus $b\in\sqrt{I}$ as $\sqrt{I}$ is invariant under taking radical. 
\par\textbf{Problem 8}\\
As we have seen $\mathfrak{p}_\mathfrak{U}$ is a prime ideal, $R_\mathfrak{U}$ is a ring. Remains to show that this is a valuation ring of $\mathfrak{k}(\mathfrak{p})$. Let $z,n\in A \backslash \mathfrak{p}_\mathfrak{U}$ and $x = {\frac z n}$. Then we have,
\begin{equation*}
    V_A = \{(\p,R)\sep v_{R,\p}(x)\leq 1\}\cup \{(\p,R)\sep v_{R,\p}(x)\geq 1\}.
\end{equation*}
By the property of ultrafilters, one of the summands belongs to $\U$. In particular, $x$ or $x^{-1}$ belong to $R_\U$. Thus it is a valuation ring.
\par Finally, to show that it is a generic limit, enough to consider basis open elements,
\begin{equation*}
    \tilde{\mathcal{R}}(f_0|f_1),
\end{equation*}
we have,
\begin{equation*}
    (\p_\U,R_\U) \in \tilde{\mathcal{R}}(f_0|f_1)\iff \tilde{\mathcal{R}}(f_0|f_1)\in\U.
\end{equation*}
$(\p_\U,R_\U)\in \tilde{\mathcal{R}}(f_0|f_1)$ if and only if 
\begin{equation*}
    f_0\not\in \p_\U,f_1\mod{\p_\U}\in f_0\mod{\p_\U}R_\U.
\end{equation*}
Thus $\tilde{\mathcal{R}}(f_0|f_1)\in\U$ as for any $(\p,R)\in\tilde{\mathcal{R}}(f_0|f_1)$, 
\begin{equation*}
    v_{R,\p}(f_1)\leq v_{R,\p}(f_0).
\end{equation*}
Thus $(\p_\U,R_\U)$ is a generic limit. They are closed under $\text{glim}$ as $\mathcal{R}(f_0|f_1,\cdots,f_n)$ are quasi-compact in $\text{Spv}A$.
\end{document}